% ============================================================================
% SYSTEM_MODELS_DRAFT.tex — polished "System Architecture and Models" (REVIEW COPY).
% Flowing rewrite organized around the shared-camera coupling. Data-independent.
% Keeps the existing TikZ system figure (fig:system) from ACODS.tex (referenced, not
% redrawn here). Notation consistent with THEORY_DRAFT.tex / RELATED_WORK_DRAFT.tex.
% Folds in: (i) the E0 empirical grounding of R(z) ~ z^2, (ii) the A_K = I - dt K
% reconciliation, (iii) forward references into the theory (Prop. A/E, Cor. 1).
% ============================================================================

\section{System Architecture and Models}
\label{sec:arch}

A UAV tracks a UGV in a common world frame $\mathcal{F}_W$. The UAV carries the
system's \emph{only} exteroceptive sensor---a single downward-facing monocular
camera---and the UGV is otherwise perception-limited, equipped only with wheel
odometry. From each camera frame the system extracts two products: the lane ahead of
the vehicle, segmented to yield a navigation waypoint, and an ArUco marker on the UGV,
detected to yield a noisy pose measurement. The pose measurement is fused with odometry
in an extended Kalman filter (EKF), producing a belief $(\mu,\Sigma)$ over the UGV
state. The single camera therefore plays a \emph{dual role}: it is at once the UGV's
navigation reference and the sole input to its state estimator (Fig.~\ref{fig:system}).

This shared-sensor design is the source of the coupling studied in this paper. Because
both products are read from the same image, they share the same imaging geometry, and
in particular the same altitude. Raising the UAV enlarges the ground footprint and so
improves navigation coverage, but it also coarsens the ground sampling distance and so
degrades localization; lowering the UAV does the reverse. A single scalar---UAV
altitude---thus governs estimation quality and navigation coverage \emph{simultaneously
and in opposite directions}. The remainder of this section makes the two sides of this
tension precise: the UGV and UAV motion models, the altitude-dependent measurement
model (the resolution side), and the coverage model (the navigation side). The control
objective that ties them together is threefold: keep the UGV within the footprint so
that $\Sigma$ stays bounded, track it robustly under bounded disturbances, and do so
while limiting the frequency of online re-optimization.

\subsection{Ground vehicle dynamics}
The UGV follows unicycle kinematics with first-order actuator lag. With state
$\mathbf{x}_g=[x_g,y_g,\theta,v,\omega]^{\!\top}$---planar position $(x_g,y_g)$,
heading $\theta$, and linear/angular rates $(v,\omega)$---and commanded velocities
$\mathbf{u}_g=[v_c,\omega_c]^{\!\top}$, the discrete-time model is
\begin{equation}
\mathbf{x}_{g,k+1}=\mathbf{f}_g(\mathbf{x}_{g,k},\mathbf{u}_{g,k})+\mathbf{w}_g,
\label{eq:ugv}
\end{equation}
where the actuator lag is set by time constants $\tau_v,\tau_\omega$ and
$\mathbf{w}_g$ is bounded process noise. The heading and rate states are retained in
the estimator because the marker yields a full planar pose, not merely a position.

\subsection{Aerial vehicle model}
Although the platform is a quadrotor, the planning and tracking layers treat it as a
velocity-controlled point mass; an inner-loop attitude/thrust controller (the Gazebo
velocity-control plugin in our study) is assumed to realize commanded velocities. With
position $\mathbf{p}=[p_x,p_y,p_z]^{\!\top}$, altitude $z:=p_z$, and velocity command
$\mathbf{u}=[v_x,v_y,v_z]^{\!\top}$, the single-integrator model is
\begin{equation}
\mathbf{p}_{k+1}=\mathbf{p}_k+\Delta t\,\mathbf{u}_k+\mathbf{w}_k,
\qquad A_d=I_3,\quad B_d=\Delta t\,I_3,
\label{eq:uav_disc}
\end{equation}
with bounded disturbance $\mathbf{w}_k\in\mathbb{W}$. All three axes are actuated, so
$v_z$ is a genuine decision variable and the altitude policy is realized in closed
loop. The linear form \eqref{eq:uav_disc} is what later makes the ancillary error
dynamics exact: with the tracking law $\mathbf{u}=\bar{\mathbf{u}}-K\mathbf{e}$ the
closed-loop error map is $A_K=A_d-B_dK=I_3-\Delta t\,K$, contractive for the deployed
gain (Sec.~\ref{sec:theory-formal}, Prop.~\ref{prop:iss}).

\subsection{Resolution model: altitude-dependent measurement noise}
\label{sec:gsd}
The accuracy of marker-based localization is governed by the camera's ground sampling
distance (GSD), the metric size of one pixel on the ground, which grows linearly with
altitude. If the image-plane detection error is approximately constant in pixels, the
metric position error grows in proportion to the GSD, and the position-measurement
covariance therefore grows quadratically with altitude. Let $z_0$ be a reference
altitude, $d=\lVert\mathbf{p}^{xy}-\mathbf{p}_g\rVert$ the horizontal UAV--UGV
separation, and $r_{\mathrm{fov}}(z)=z\tan(\phi/2)$ the footprint radius for camera
field of view $\phi$. We model the measurement covariance as
\begin{equation}
R(z)=R_0\!\left(\frac{z}{z_0}\right)^{\!2}
\!\left(1+\left(\frac{d}{r_{\mathrm{fov}}}\right)^{\!2}\right),
\label{eq:Rgsd}
\end{equation}
where $R_0$ is the baseline covariance at $z_0$ and the second factor mildly inflates
the noise as the target approaches the footprint edge. Equation~\eqref{eq:Rgsd} is
strictly increasing in $z$; this monotonicity is the resolution side of the coupling
and the hypothesis of the optimal-altitude result (Prop.~\ref{prop:alt}). Importantly,
the law is not merely postulated: in static-hover trials the measured localization
error grows linearly with altitude (median error $\approx 1.03+2.24\,z$\,mm,
$R^2=0.994$), confirming $R\propto z^2$ (Sec.~\ref{sec:results}, Fig.~\ref{fig:envelope}).
The \emph{same} model \eqref{eq:Rgsd} is used by both the EKF and the belief-space
planner in all reported experiments, so the planner propagates exactly the uncertainty
the estimator experiences.

\subsection{Coverage model: the segmentation floor}
\label{sec:coverage}
Lane segmentation seeds a waypoint a lookahead distance ahead of the vehicle, and that
waypoint must fall inside the footprint or the reference is lost. Writing the required
lookahead as $d_{\mathrm{look}}(\kappa)$---which grows with local path curvature
$\kappa$, since a curving lane leaves the footprint sooner---we summarize navigability
by the non-dimensional \emph{coverage factor}
\begin{equation}
\eta(z,\kappa)=\frac{r_{\mathrm{fov}}(z)}{d_{\mathrm{look}}(\kappa)}
             =\frac{z\tan(\phi/2)}{d_{\mathrm{look}}(\kappa)},
\label{eq:eta}
\end{equation}
with segmentation feasible for $\eta\gtrsim1$, i.e.\ $z\ge
z_{\mathrm{floor}}(\kappa)=d_{\mathrm{look}}(\kappa)/\tan(\phi/2)$. This is the
navigation side of the coupling, and it opposes the resolution side directly: the
estimator prefers \emph{low} altitude for resolution \eqref{eq:Rgsd}, while
\eqref{eq:eta} demands \emph{sufficient} altitude for coverage, and the floor
$z_{\mathrm{floor}}(\kappa)$ rises exactly where the path curves. The two models
\eqref{eq:Rgsd} and \eqref{eq:eta} therefore frame a single question---how low may the
UAV fly without losing the lane---whose answer is the estimation-optimal, curvature-
scheduled altitude derived next (Prop.~\ref{prop:alt}, Cor.~\ref{cor:climb}) and
realized by the aerial planner of Sec.~\ref{sec:method}.
